\documentclass[11pt]{article}
\usepackage[margin=1in]{geometry}
\usepackage{amsmath,amssymb,amsthm}
\usepackage[hidelinks]{hyperref}

\newtheorem{theorem}{Theorem}
\newtheorem{lemma}{Lemma}
\newtheorem{corollary}{Corollary}

\newcommand{\eps}{\varepsilon}
\newcommand{\Span}{\operatorname{span}}
\newcommand{\osc}{\operatorname{osc}}

\title{A Hilbert-Operator Partial Result for Hanson Question 4.3}
\author{fa\_banach\_001, agent\_lane\_19}
\date{2026-06-16}

\begin{document}
\maketitle

\section*{Source Question}

James Hanson asks in \emph{Indiscernible Subspaces and Minimal Wide Types},
arXiv:2004.03062, Question 4.3:
\begin{quote}
Does (model theoretic) stability imply oscillation stability? That is to say,
if \(T\) is a stable Banach theory, is every unary formula oscillation stable
on models of \(T\)?
\end{quote}
The arXiv record notes that this preprint was absorbed into arXiv:2011.00610,
\emph{Strongly Minimal Sets and Categoricity in Continuous Logic}; the corrected
absorbing paper keeps the same question in its Banach-theory section.

\section*{Result}

This note proves a positive subcase for quantifier-free formulas in Hilbert
structures expanded by bounded linear operators in the standard many-sorted
operator-linear Hilbert/Banach language.  The result does not use
model-theoretic stability.  Its relevance is that many natural examples near
Hanson's question are stable Hilbert structures with distinguished linear
operators and quantifier elimination.

\section*{Proof Intuition}

A unary quantifier-free Hilbert-operator formula only uses finitely many atoms
of the form \(\|Sx+b\|\), where \(S\) is a bounded operator obtained from the
named operators by finitely many algebraic operations and \(b\) is a parameter
vector.  On an infinite-dimensional subspace, the finitely many linear terms
\(\langle Sx,b\rangle\) can be killed by passing to a finite-codimensional
subspace.  What remains is a finite list of quadratic forms
\(\langle S^*Sx,x\rangle\).  A diagonal orthonormal-subsequence argument
produces an infinite-dimensional subspace on which all those quadratic forms
are simultaneously almost constant.  Uniform continuity of the continuous
connectives then gives small oscillation for the whole formula.

\section*{A Compression Lemma}

\begin{lemma}[simultaneous quadratic flattening]\label{lem:flatten}
Let \(H\) be an infinite-dimensional real or complex Hilbert space, let
\(A_1,\ldots,A_m\in B(H)\) be bounded self-adjoint operators, and let
\(\eta>0\).  Then there are an infinite-dimensional closed subspace \(Z\subset H\)
and scalars \(\lambda_1,\ldots,\lambda_m\in\mathbb{R}\) such that, for every
unit vector \(z\in Z\),
\[
  \left|\langle A_j z,z\rangle-\lambda_j\right|<\eta
  \qquad (1\leq j\leq m).
\]
\end{lemma}

\begin{proof}
Choose an orthonormal sequence \((u_n)\) in \(H\).  By compactness of
\(\prod_{j=1}^m[-\|A_j\|,\|A_j\|]\), after passing to a subsequence we may
assume that \(\langle A_j u_n,u_n\rangle\to\lambda_j\) for every \(j\).

We now pass to a further subsequence \((e_n)\) recursively.  Since every
orthonormal sequence is weakly null, once \(e_1,\ldots,e_{n-1}\) have been
chosen, a sufficiently far tail vector \(u_k\) satisfies
\[
  |\langle A_j u_k,u_k\rangle-\lambda_j|<\eta/4
\]
for all \(j\), and also
\[
  |\langle A_j u_k,e_r\rangle|
     < \frac{\eta}{16m}2^{-n-r}
     \qquad (1\leq j\leq m,\; r<n).
\]
Relabel the resulting subsequence as \((e_n)\), and let
\(Z=\overline{\Span}\{e_n:n\geq1\}\).

Fix \(j\) and write a unit vector \(z\in Z\) as \(z=\sum_n c_n e_n\).  Then
\[
\begin{aligned}
\left|\langle A_jz,z\rangle-\lambda_j\right|
&\leq
\sum_n |c_n|^2
   \left|\langle A_j e_n,e_n\rangle-\lambda_j\right|  \\
&\quad+
\sum_{n\neq r}|c_n|\,|c_r|\,|\langle A_j e_n,e_r\rangle|.
\end{aligned}
\]
The diagonal part is \(<\eta/4\).  For the off-diagonal part, use
\((\sum_n2^{-n}|c_n|)^2\leq\sum_n4^{-n}<1\).  The recursive bounds, together
with self-adjointness for the reversed pairs, make this part \(<\eta/(16m)\).
Hence the whole expression is \(<\eta\).
\end{proof}

\section*{The Partial Theorem}

\begin{theorem}\label{thm:qf}
Let \(H\) be an infinite-dimensional Hilbert space in the standard many-sorted
operator-linear Hilbert/Banach language, expanded by any collection of bounded
linear operator symbols.  Let \(\varphi(x,\bar a)\) be a quantifier-free unary
formula with parameters \(\bar a\), whose one-variable atoms are norms of
operator-linear terms.  For every infinite-dimensional closed subspace
\(Y\subseteq H\) and every \(\eps>0\), there exists an infinite-dimensional
closed subspace \(Z\subseteq Y\) such that
\[
  \sup\{|\varphi(u,\bar a)-\varphi(v,\bar a)|:\ u,v\in S_Z\}<\eps .
\]
\end{theorem}

\begin{proof}
Only finitely many atomic predicates occur in \(\varphi\).  By the
operator-linear language hypothesis, each unary term is of the form \(T x+c\),
where \(T\) is a bounded operator obtained as a finite linear combination of
finite words in the named operator symbols, and \(c\) is a parameter vector.
Hence each unary atomic predicate appearing in \(\varphi\) has the form
\[
  x\longmapsto \|Sx+b\|
\]
for some bounded operator \(S\) and vector \(b\).

Let \(\|S_1x+b_1\|,\ldots,\|S_Nx+b_N\|\) list these atoms.  By uniform
continuity of the continuous connective defining \(\varphi\), choose
\(\delta>0\) so small that changing every atomic value by at most \(\delta\)
changes \(\varphi\) by less than \(\eps\).

Inside \(Y\), remove the finitely many linear terms by setting
\[
  Y_0
  =
  Y\cap \bigcap_{\ell=1}^N
  \ker\bigl(z\mapsto \langle S_\ell z,b_\ell\rangle\bigr).
\]
This is still infinite-dimensional.  For \(z\in Y_0\) with \(\|z\|=1\),
\[
  \|S_\ell z+b_\ell\|^2
  =
  \langle S_\ell^*S_\ell z,z\rangle+\|b_\ell\|^2 .
\]

Apply Lemma~\ref{lem:flatten} in \(Y_0\) to the finite self-adjoint family
\(S_\ell^*S_\ell\) restricted to \(Y_0\), with tolerance chosen small enough
that the corresponding square-rooted norms vary by at most \(\delta\).  The
result is an infinite-dimensional closed subspace \(Z\subseteq Y_0\) on which
every atom \(\|S_\ell x+b_\ell\|\) varies by at most \(\delta\) over \(S_Z\).
Therefore \(\varphi(x,\bar a)\) varies by less than \(\eps\) over \(S_Z\), as
required.
\end{proof}

\section*{Transfer to Quantifier-Elimination Classes}

\begin{corollary}[QE transfer]\label{cor:qe-transfer}
Let \(T\) be a complete theory of infinite-dimensional Hilbert spaces in a
language obtained by adding bounded linear operator symbols to the standard
Hilbert/Banach language.  Suppose that every unary formula of \(T\) is, over
\(T\), uniformly approximable by quantifier-free unary formulas whose atoms are
norms of operator-linear terms.  Then every unary formula of \(T\) is
oscillation stable on every model of \(T\).
\end{corollary}

\begin{proof}
Let \(\varphi(x,\bar a)\) be a unary formula, \(Y\) an infinite-dimensional
closed subspace of a model of \(T\), and \(\eps>0\).  Choose a quantifier-free
\(\psi(x,\bar a)\) of the kind handled in Theorem~\ref{thm:qf} such that
\(|\varphi-\psi|<\eps/3\) on the relevant sort.  Theorem~\ref{thm:qf} gives an
infinite-dimensional closed \(Z\subseteq Y\) on whose unit sphere \(\psi\) has
oscillation \(<\eps/3\).  Then \(\varphi\) has oscillation \(<\eps\) on
\(S_Z\).
\end{proof}

\begin{corollary}[named stable Hilbert-operator classes]\label{cor:named}
Hanson's Question~4.3 has a positive answer for the following stable
Hilbert-operator theories.
\begin{enumerate}
\item Completions of the theory of Hilbert spaces expanded by a bounded normal
operator \(T\), in the language with \(T\) and \(T^*\).
\item The self-adjoint and unitary one-operator subcases in the original
language with just the named operator.
\item Hilbert spaces expanded by a unitary representation of a finite group.
\end{enumerate}
\end{corollary}

\begin{proof}
For (1), Berenstein--Cuervo Ovalle--Goldbring prove quantifier elimination for
the completions after adding \(T^*\) (Theorem 4.1 of
arXiv:2507.21894) and prove that the completions are stable, indeed
superstable (Theorems 6.2 and 6.3 of arXiv:2507.21894).  Their quantifier-free
terms are built from \(T,T^*\), hence are operator-linear terms covered by
Corollary~\ref{cor:qe-transfer}.

The same paper proves quantifier elimination in the original language for the
self-adjoint and unitary cases (Theorem 4.3 of arXiv:2507.21894), so (2)
follows in the same way.  For (3), Berenstein--Perez prove quantifier
elimination for Hilbert spaces expanded by a unitary representation of a finite
group (Proposition 3.6 of arXiv:2409.03923) and \(\aleph_0\)-stability
(Theorem 3.8 of arXiv:2409.03923).  Again the terms are finite linear
combinations of named unitary actions, so Corollary~\ref{cor:qe-transfer}
applies.
\end{proof}

\section*{Limitations}

Theorem~\ref{thm:qf} is not a full answer to Hanson Question 4.3.  It applies
to quantifier-free unary formulas in the standard operator-linear
Hilbert/Banach language.  Corollary~\ref{cor:named} gives full positive
answers for several named stable Hilbert-operator theories, but it still does
not cover arbitrary stable Banach theories or stable expansions by general
unary predicates.

\section*{Novelty Check}

The run indexes were searched for arXiv:2004.03062, \emph{Indiscernible
Subspaces}, minimal wide types, and oscillation stability.  Web searches for
the exact phrase ``Does model theoretic stability imply oscillation stability''
and close variants found the source and the absorbing paper, but no explicit
later answer.  Further searches found the normal-operator and finite-group
Hilbert expansion papers cited below; these provide named stable QE classes to
which the transfer corollary applies.

\section*{References}

\begin{itemize}
\item James Hanson, \emph{Indiscernible Subspaces and Minimal Wide Types},
  arXiv:2004.03062.
\item James Hanson, \emph{Strongly Minimal Sets and Categoricity in Continuous
  Logic}, arXiv:2011.00610.
\item Alexander Berenstein, Nicol\'as Cuervo Ovalle, and Isaac Goldbring,
  \emph{Model theory of Hilbert spaces expanded by normal operators},
  arXiv:2507.21894.
\item Alexander Berenstein and Juan Manuel Perez,
  \emph{Model Theory of Hilbert Spaces with a Discrete Group Action},
  arXiv:2409.03923.
\end{itemize}

\end{document}
